\section{Exact constants for the single-cube criterion}\label{app:sc-certificates}

The estimates below certify only the finite constants used in
Sections~\ref{sc:sec:psts}--\ref{sc:sec:continuum}. They do not estimate a
Wilson partition function at weak coupling. The accompanying
\texttt{anc/single\_cube\_constants.py} uses rational arithmetic to
reproduce every comparison in this appendix.

For $y\ge0$, let $S_m(y)=\sum_{k=0}^m y^k/k!$. If $y<m+2$, the exponential
series and the geometric bound on its tail give
\begin{equation}\label{sc:eq:exp-certificate}
S_m(y)\le e^y\le U_m(y):=S_m(y)+
\frac{y^{m+1}}{(m+1)!}\left(1-\frac{y}{m+2}\right)^{-1}.
\end{equation}
Indeed the ratio of any two consecutive tail terms is at most $y/(m+2)$.
Also $e^y-1\le y/(1-y)$ for $0\le y<1$, because $e^{-y}\ge1-y$.
These facts reduce the needed comparisons to the following finite
rational inequalities:
\begin{align*}
&\frac{2}{8192}-\frac{3}{8192^2}-\frac1{4097}
=\frac{4093}{274945015808}>0;\\
&64\cdot252=16128<16129=127^2,
\qquad \frac{252}{251}\le\frac{127}{126};\\
&4160\left(U_4\left(\frac{64}{4160}\right)-1\right)^2
<\frac{99991}{100000}<1;\\
&4159\left(S_4\left(\frac{64}{4159}\right)-1\right)^2
>\frac{50007}{50000}>1;\\
&U_8\left(\frac{677}{1000}\right)<\frac{246}{125}<\frac{63}{32}.
\end{align*}
All denominators in these expressions are positive; clearing them gives
integer comparisons. The first two lines justify purity conversion and
the inductive envelope. The next two bracket the unique positive fixed
point of $f(x)=(e^{64x}-1)^2$ by $1/4160$ and $1/4159$.
The last line implies $e^{0.677}<63/32$, hence
$\log(63/32)>0.677$. No rounded floating-point output is used in any of
these deductions.

\begin{remark}[zero eigenvalues]
Every gap statement in these sections uses $-\log0=+\infty$.
When all excited eigenvalues vanish, $\mathcal Z_s(T)=1$ for all $T\ge2$.
The resulting finite-temperature estimates remain valid, but this
rank-one example is a reminder that an arbitrarily large spectral gap
alone does not prove that a continuum theory has any nonvacuum states.
\end{remark}
